\section{Differentiation in \(\RR^n\)}

In this section, we assume the norm \(\norm{\cdot}\) on \(\RR^n\) to be the Euclidean norm.

\subsection{Limits and Differentiability}

First, we need the following definition of limit.

\begin{definition}[Limit in \(\RR^n\)]
    Let \(E \subseteq \RR^n\) and \(\functiondomain{f}{E}{\RR^n}\). Let \(a\) be a limit point of the set \(E\), and let \(b \in \RR^n\). We say the \term{limit of \(x \to a\) of \(f\pqty{x}\) is \(b\)}, denoted \(\lim_{x \to a} f\pqty{x} = b\), if for every \(\epsilon > 0\), there exists \(\delta > 0\) such that for all \(x \in E\) with \(0 < \norm{x - a} < \delta\), we have \(\norm{f\pqty{x} - b} < \epsilon\), i.e.,
    \[
        \forall \epsilon > 0, \exists \delta > 0, \forall x \in E: 0 < \norm{x - a} < \delta \implies \norm{f\pqty{x} - b} < \epsilon.
    \]
\end{definition}

\begin{remark}
    We do not assume that \(a \in E\) (i.e. whether \(f\) is defined at \(E\), or if so, the value of \(f\) at \(E\), does not matter).
\end{remark}

How should we define differentiability in \(\RR^n\), say a real-value function \(\functiondomain{f}{\RR^n}{\RR}\), or perhaps, \(\functiondomain{f}{\RR^n}{\RR^m}\) at a point \(x \in \RR^n\)? We want to come up with a definition that is consistent with the one-dimensional case of \(n = 1\), a function \(\functiondomain{f}{\RR}{\RR}\) is \term{differentiable} at \(x \in \RR\) if the limit
\[
    D f\pqty{a} = f'\pqty{a} = \lim_{h \to 0} \frac{f\pqty{a + h} - f\pqty{a}}{h}.
\]

This limit makes sense for \(f\) defined on \(\RR^n\) as well, but what if the function not real-valued, but real-vector-valued? We might naturally define a limit
\[
    \lim_{\vb{h} \to \vb{0}} \frac{f\pqty{a + \vb{h}} - f\pqty{a}}{\norm{\vb{h}}},
\]
where we divide the norm of a vector instead of a real number. This immediately gives inconsistency in the one-dimensional case! If we apply this definition to \(n = 1\), we have \(\norm{h} = \abs{h}\), and the limit becomes
\[
    \lim_{h \to 0} \frac{f\pqty{a + h} - f\pqty{a}}{\abs{h}},
\]
which is very different. For example, take \(f\pqty{x} = \abs{x}\) at \(a = 0\), then this new limit exists and is equal to \(1\), but clearly in the traditional definition, \(f\) is not differentiable at \(0\). On the other hand, take \(f\pqty{x} = x\) at \(x = 0\), then the new limit does not exist, but clearly \(f\) is differentiable at \(0\) with derivative \(1\).

But note that in the \(n = 1\) case, the value of the derivative \(f'\pqty{x}\) is equivalent to saying
\[
    \lim_{h \to 0} \frac{f\pqty{x + h} - f\pqty{x} - h f'\pqty{x}}{h} = 0,
\]
and since this limit takes a value of \(0\), there is nothing stopping us from taking the absolute value! So in other words, \(f\) is differentiable if and only if
\[
    \lim_{\vb{h} \to \vb{0}} \frac{f\pqty{x + \vb{h}} - f\pqty{x} - A\vb{h}}{\abs{\vb{h}}} = 0.
\]
for some real number \(A \in \RR\), which would be the value of the derivative  But now, this definition can be simply generalised to a higher dimensional case.

\begin{definition}[Differentiability in \(\RR^n\)]
    Let \(U \subseteq \RR^n\) be an open set, and let \(\functiondomain{f}{U}{\RR^m}\), and take \(a \in U\). We say that \(f\) is \term{differentiable at \(a\)} if there exists a linear map \(A \in \linearmap\pqty{\RR^n, \RR^m}\) such that
    \[
        \lim_{\vb{h} \to \vb{0}} \frac{f\pqty{a + \vb{h}} - f\pqty{a} - A\vb{h}}{\norm{\vb{h}}} = 0.
    \]

    We say \(A\) is the \term{derivative of \(f\) at \(a\)}, and we denote it by \(Df\pqty{a} = A\).
\end{definition}

\begin{remark}
    We require \(U\) to be open, so that when \(a \in U\), there is \(r > 0\) such that \(\oball{x}{r} \subseteq U\), so \(a + h \in U\) for \(h\) with sufficiently small norm.
\end{remark}

\begin{proposition}[Properties of Derivatives]
    Let \(U \subseteq \RR^n\) be open, \(\functiondomain{f}{U}{\RR^m}\), and let \(a \in U\). We have the following.
    \begin{enumerate}[label=(\roman*)]
        \item If \(f\) is differentiable at \(a\), then the derivative \(Df\pqty{a}\) is unique.
        \item If \(f\) is differentiable at \(a\), then \(f\) is continuous at \(a\).
        \item Write \(f\) in components \(f = \pqty{f_1, \dots, f_m}\), where \(\functiondomain{f_j}{U}{\RR}\) for \(j = 1, \dots, m\). Then \(f\) is differentiable at \(a\) if and only if each component \(f_i\) is differentiable at \(a\), and in this case, we have
        \[
            f \pqty{a} = \pmqty{D f_1\pqty{a} \\ D f_2 \pqty{a} \\ \vdots \\ Df_m\pqty{a}}.
        \]
    \end{enumerate}
\end{proposition}

\begin{proof}
    \begin{enumerate}[label=(\roman*)]
        \item If \(A\) and \(A'\) are different derivatives of \(f\) at \(a\), then we have
        \[
            \lim_{\vb{h} \to \vb{0}} \frac{f\pqty{a + \vb{h}} - f\pqty{a} - A\vb{h}}{\norm{\vb{h}}} = 0.
        \]

        For any \(\vb{u} \in \RR^n \setminus \set{0}\), we take \(\vb{h} = t\vb{u}\), then
        \[
            \lim_{t \to 0} \frac{t A \vb{u} - tA'\vb{u}}{\norm{t\vb{u}}} = \lim_{t \to 0} \frac{t \pqty{A - A'} \vb{u}}{\abs{t}\norm{\vb{u}}} = \lim_{t \to 0} \frac{\pqty{A - A'}\vb{u}}{\norm{\vb{u}}} = 0,
        \]
        and hence \(A\vb{u} = A'\vb{u}\) for all \(\vb{u} \in \RR^n\), and hence \(A = A'\).

        \item We have
        \[
            \lim_{\vb{h} \to \vb{0}} \pqty{f\pqty{a + \vb{h}} - f\pqty{a} - A\vb{h}} = 0,
        \]
        and hence
        \[
            \lim_{\vb{h} \to \vb{0}} \pqty{f\pqty{a + \vb{h}} - f\pqty{a}} = \lim_{\vb{h} \to \vb{0}} A\vb{h} = 0,
        \]
        and hence \(f\) is continuous at \(a\).

        \item This follows from the fact that for any \(A = \pmqty{A_1 \\ A_2 \\ \vdots \\ A_m} \in \linearmap\pqty{\RR^n, \RR^m}\), we have \(A_j \in \linearmap\pqty{\RR^n, \RR}\) for \(j = 1, \dots, m\), and by inspecting components,
        \[
            \lim_{\vb{h} \to \vb{0}} \frac{f\pqty{a + \vb{h}} - f\pqty{a} - A\vb{h}}{\norm{\vb{h}}} = 0 \iff \lim_{\vb{h} \to \vb{0}} \frac{f_j\pqty{a + \vb{h}} - f_j\pqty{a} - A_j\vb{h}}{\norm{\vb{h}}} = 0.
        \]
    \end{enumerate}
\end{proof}

\subsection{Directional Derivatives and Partial Derivatives}

\begin{definition}[Directional Derivative]
    Let \(\functiondomain{f}{U}{\RR^m}\) where \(U \subseteq \RR^n\) is open. Let \(\vb{u} \in \RR^n\) be fixed and \(a \in U\). The \term{directional derivative} of \(f\) at \(a\) in the direction \(\vb{u}\) is
    \[
        D_{\vb{u}} f\pqty{a} = \lim_{t \to 0} \frac{f\pqty{a + t\vb{u}} - f\pqty{a}}{t},
    \]
\end{definition}

\begin{proposition}[Computation of Directional Derivatives]
    If \(f\) is differentiable at \(x = a\), then for any \(\vb{u} \in \RR^n\), the directional derivative \(D_{\vb{u}} f\pqty{a}\) exists and is given by
    \[
        D_{\vb{u}} f\pqty{a} = Df\pqty{a}\vb{u}.
    \]

    In particular, the map \(\vb{u} \mapsto D_{\vb{u}} f\pqty{a}\) is linear if \(f\) is differentiable at \(a\).
\end{proposition}

\begin{proof}
    Since \(f\) is differentiable at \(a\), taking \(h = t\vb{u}\) for \(t \in \RR\), we have
    \[
        \lim_{t \to 0} \frac{f\pqty{a + t\vb{u}} - f\pqty{a} - Df\pqty{a} t\vb{u}}{\norm{t\vb{u}}} = 0,
    \]
    and hence
    \[
        \lim_{t \to 0} \frac{f\pqty{a + t\vb{u}} - f\pqty{a}}{t} = Df\pqty{a}\vb{u}.
    \]
\end{proof}

In the special case of \(m = 1\) under the standard basis, we may define a special case of the directional derivative.

\begin{definition}[Partial Derivative]
    Let \(\functiondomain{f}{U}{\RR}\) where \(U \subseteq \RR^n\) is open, and let
    \[
        \vb{e}_j = \pmqty{0 \\ \vdots \\ 1 \\ \vdots \\ 0} \in \RR^n
    \]
    be the standard basis vector in \(\RR^n\) with a \(1\) in the \(j\)-th position and \(0\) elsewhere. We denote the \term{\(j\)-th partial derivative} of \(f\) at \(a\) by
    \[
        \pdv{f}{x_j}\pqty{a} = D_j f\pqty{a} = D_{\vb{e}_j} f\pqty{a} = \lim_{t \to 0} \frac{f\pqty{a + t\vb{e}_j} - f\pqty{a}}{t}.
    \]
\end{definition}